\documentclass[12pt]{article}
\usepackage[utf8]{inputenc}
\usepackage[T1]{fontenc}
\usepackage{hyperref}
\title{Chesu Language Resources for K-12 Curriculum Development}
\author{Compiled by Claude Sonnet 5.5 (claude-sonnet-5-5) for Andr\'as Kornai}
\date{}
\begin{document}
\maketitle

\section{Introduction}
Chesu (ISO 639-3: ych) is a Loloish (Sino-Tibetan) language of Yunnan, China. It is spoken in southern Shuangbai County, northern Xinping County and Eshan County. Wikipedia cites Ethnologue for about 3,300 speakers (2007) out of an ethnic population of about 6,600. Wikipedia reports Bradley (2007) as classifying it as Nasoid, closely related to Eastern Yi Nasu. Officially the Chesu are counted as part of the Yi nationality, but they do not consider themselves Yi. Some Chesu no longer speak the language and use Chinese instead.

\section{Primary Educational Resources}
The search found no primers, textbooks or school programmes for Chesu.

\section{Linguistic and Reference Materials}
\begin{itemize}
\item Wikipedia overview: \url{https://en.wikipedia.org/wiki/Chesu_language}
\item Endangered Languages Project entry: \url{https://endangeredlanguages.com/lang/7251}
\item Verbix endangered-languages page: \url{https://docs.verbix.com/EndangeredLanguages/Chesu}
\end{itemize}
The search found no grammar or dictionary of Chesu.

\section{Religious and Community Resources}
\begin{itemize}
\item Joshua Project people-group profile: \url{https://joshuaproject.net/people_groups/18429}
\item Joshua Project language page: \url{https://joshuaproject.net/languages/ych}
\item Asia Harvest people-group profile (PDF): \url{https://people-groups.asiaharvest.org/people-group-profiles/yunnan/chesu.pdf}
\end{itemize}
I did not read these pages. I cannot say whether they hold any Chesu-language content.

\section{Digital and Audio Resources}
The search found no audio or digital resources specific to Chesu.

\section{Curriculum Development Context}
Chesu appears to have very little documentation or teaching material. A curriculum would need documentation first, with Mandarin as the school language. The Chesu are administratively counted as Yi, and Yi-language educational resources exist elsewhere. The Chesu may not regard them as relevant to their own language.

\end{document}